% Worked Examples: Concept 1.3.1 - Basic Definitions
\documentclass[11pt,a4paper]{article}
\usepackage{worked-examples-template}

\title{\textbf{Worked Examples}\\[0.5em]
\large Concept 1.3.1: Basic Definitions\\[0.3em]
\normalsize \coursesubtitle{} -- \coursetitle}
\author{Assoc.\ Prof.\ William Robertson \and Dr Sean McGowan}
\date{\today}

\begin{document}

\maketitle
\tableofcontents
\newpage

\section{Concept 1.3.1: Linear Systems}

\subsection{Example 1: Linearity and the Superposition Principle}

\paragraph{Problem:} Consider a linear time-invariant (LTI) system described by:
\begin{align}
\frac{d^2y}{dt^2} + 3\frac{dy}{dt} + 2y = u(t)
\end{align}

(a) Verify that the system is linear by checking if it satisfies superposition.

(b) Find the responses $y_1(t)$ and $y_2(t)$ to inputs $u_1(t) = e^{-t}$ and $u_2(t) = e^{-2t}$ respectively, with zero initial conditions.

(c) Using superposition, determine the response to $u_3(t) = 3e^{-t} + 2e^{-2t}$.

(d) Explain why superposition does not apply if initial conditions are non-zero.

\paragraph{Solution:}

\textbf{Step 1: Linearity verification}

A system is linear if it satisfies:
\begin{itemize}
\item \textbf{Additivity:} If $u_1 \to y_1$ and $u_2 \to y_2$, then $(u_1 + u_2) \to (y_1 + y_2)$
\item \textbf{Homogeneity:} If $u \to y$, then $\alpha u \to \alpha y$ for any scalar $\alpha$
\end{itemize}

For our differential equation, let $u_1 \to y_1$ and $u_2 \to y_2$:
\begin{align}
\frac{d^2y_1}{dt^2} + 3\frac{dy_1}{dt} + 2y_1 &= u_1 \\
\frac{d^2y_2}{dt^2} + 3\frac{dy_2}{dt} + 2y_2 &= u_2
\end{align}

Consider $y_3 = \alpha y_1 + \beta y_2$:
\begin{align}
\frac{d^2y_3}{dt^2} + 3\frac{dy_3}{dt} + 2y_3 &= \frac{d^2(\alpha y_1 + \beta y_2)}{dt^2} + 3\frac{d(\alpha y_1 + \beta y_2)}{dt} + 2(\alpha y_1 + \beta y_2) \\
&= \alpha\left(\frac{d^2y_1}{dt^2} + 3\frac{dy_1}{dt} + 2y_1\right) + \beta\left(\frac{d^2y_2}{dt^2} + 3\frac{dy_2}{dt} + 2y_2\right) \\
&= \alpha u_1 + \beta u_2
\end{align}

Therefore, the system satisfies superposition ✓

\textbf{Step 2: Response to $u_1(t) = e^{-t}$}

The transfer function is:
\begin{align}
G(s) = \frac{Y(s)}{U(s)} = \frac{1}{s^2 + 3s + 2} = \frac{1}{(s+1)(s+2)}
\end{align}

For $u_1(t) = e^{-t}$, $U_1(s) = \frac{1}{s+1}$:
\begin{align}
Y_1(s) = G(s)U_1(s) = \frac{1}{(s+1)(s+2)} \cdot \frac{1}{s+1} = \frac{1}{(s+1)^2(s+2)}
\end{align}

Partial fraction decomposition:
\begin{align}
\frac{1}{(s+1)^2(s+2)} = \frac{A}{s+1} + \frac{B}{(s+1)^2} + \frac{C}{s+2}
\end{align}

Multiply by $(s+1)^2(s+2)$:
\begin{align}
1 = A(s+1)(s+2) + B(s+2) + C(s+1)^2
\end{align}

Setting $s = -1$: $1 = B(1) \Rightarrow B = 1$

Setting $s = -2$: $1 = C(1) \Rightarrow C = 1$

Setting $s = 0$: $1 = 2A + 2B + C = 2A + 2 + 1 \Rightarrow A = -1$

Therefore:
\begin{align}
Y_1(s) = \frac{-1}{s+1} + \frac{1}{(s+1)^2} + \frac{1}{s+2}
\end{align}

Taking inverse Laplace transform:
\begin{align}
y_1(t) = -e^{-t} + te^{-t} + e^{-2t} = (t-1)e^{-t} + e^{-2t}
\end{align}

\textbf{Step 3: Response to $u_2(t) = e^{-2t}$}

For $u_2(t) = e^{-2t}$, $U_2(s) = \frac{1}{s+2}$:
\begin{align}
Y_2(s) = \frac{1}{(s+1)(s+2)} \cdot \frac{1}{s+2} = \frac{1}{(s+1)(s+2)^2}
\end{align}

Partial fractions:
\begin{align}
\frac{1}{(s+1)(s+2)^2} = \frac{A}{s+1} + \frac{B}{s+2} + \frac{C}{(s+2)^2}
\end{align}

Setting $s = -1$: $1 = A(1) \Rightarrow A = 1$

Setting $s = -2$: $1 = C(-1) \Rightarrow C = -1$

Setting $s = 0$: $1 = 2A + \frac{B}{2} + \frac{C}{4} = 2 + \frac{B}{2} - \frac{1}{4} \Rightarrow B = -\frac{3}{2}$

Wait, let me recalculate. Multiply by $(s+1)(s+2)^2$:
\begin{align}
1 = A(s+2)^2 + B(s+1)(s+2) + C(s+1)
\end{align}

$s = -2$: $1 = C(-1) \Rightarrow C = -1$

$s = -1$: $1 = A(1) \Rightarrow A = 1$

$s = 0$: $1 = 4A + 2B + C = 4 + 2B - 1 \Rightarrow B = -1$

Therefore:
\begin{align}
Y_2(s) = \frac{1}{s+1} - \frac{1}{s+2} - \frac{1}{(s+2)^2}
\end{align}

\begin{align}
y_2(t) = e^{-t} - e^{-2t} - te^{-2t} = e^{-t} - (1+t)e^{-2t}
\end{align}

\textbf{Step 4: Response to $u_3(t) = 3e^{-t} + 2e^{-2t}$ using superposition}

By superposition:
\begin{align}
y_3(t) &= 3y_1(t) + 2y_2(t) \\
&= 3[(t-1)e^{-t} + e^{-2t}] + 2[e^{-t} - (1+t)e^{-2t}] \\
&= 3(t-1)e^{-t} + 3e^{-2t} + 2e^{-t} - 2(1+t)e^{-2t} \\
&= [3t - 3 + 2]e^{-t} + [3 - 2 - 2t]e^{-2t} \\
&= (3t - 1)e^{-t} + (1 - 2t)e^{-2t}
\end{align}

\textbf{Step 5: Non-zero initial conditions}

When initial conditions are non-zero, the total response consists of:
\begin{align}
y_{\text{total}} = y_{\text{ZIR}} + y_{\text{ZSR}}
\end{align}

where:
\begin{itemize}
\item $y_{\text{ZIR}}$ = zero-input response (response to initial conditions with $u = 0$)
\item $y_{\text{ZSR}}$ = zero-state response (response to input with zero initial conditions)
\end{itemize}

Superposition applies only to the zero-state response. If we combine two inputs with different initial conditions, the initial condition contribution does not follow superposition because it depends on the specific initial state, not the input.

\textbf{Key Insights:}
\begin{itemize}
\item Superposition is the defining property of linear systems
\item For LTI systems, outputs to scaled and summed inputs can be scaled and summed
\item This principle enables powerful analysis techniques (Laplace transform, frequency response)
\item Initial conditions must be zero for superposition to apply to the total response
\item Nonlinear terms (e.g., $y^2$, $\sin y$) violate superposition
\item Linearity makes complex problems tractable by decomposition
\end{itemize}

\subsection{Example 2: Time Invariance Property}

\paragraph{Problem:} A system has input-output relationship:
\begin{align}
y(t) = \int_0^t e^{-(t-\tau)} u(\tau) d\tau
\end{align}

(a) Determine if the system is time-invariant.

(b) Show that for input $u(t)$, the output is the convolution $y(t) = h(t) * u(t)$ where $h(t) = e^{-t}$ for $t \geq 0$.

(c) Verify the time-invariance property for a specific input.

\paragraph{Solution:}

\textbf{Step 1: Time invariance test}

A system is time-invariant if a time shift in the input causes the same time shift in the output:
\begin{align}
u(t) \to y(t) \quad \text{implies} \quad u(t - t_0) \to y(t - t_0)
\end{align}

Given $u(t) \to y(t) = \int_0^t e^{-(t-\tau)} u(\tau) d\tau$

For shifted input $u(t - t_0)$:
\begin{align}
y'(t) = \int_0^t e^{-(t-\tau)} u(\tau - t_0) d\tau
\end{align}

Let $\sigma = \tau - t_0$, so $\tau = \sigma + t_0$ and $d\tau = d\sigma$:
\begin{align}
y'(t) = \int_{-t_0}^{t-t_0} e^{-(t-\sigma-t_0)} u(\sigma) d\sigma
\end{align}

Hmm, this doesn't match. Let me reconsider. The issue is the lower limit of integration (0) is fixed, not relative to the input.

Actually, for a proper LTI system, the integral should be from $-\infty$ to $t$:
\begin{align}
y(t) = \int_{-\infty}^t h(t-\tau) u(\tau) d\tau
\end{align}

If the system is causal (zero response before $t=0$), we can write:
\begin{align}
y(t) = \int_0^t e^{-(t-\tau)} u(\tau) d\tau \quad \text{for } t \geq 0, u(\tau) = 0 \text{ for } \tau < 0
\end{align}

Now check time invariance. The shifted output should be:
\begin{align}
y(t - t_0) = \int_0^{t-t_0} e^{-(t-t_0-\tau)} u(\tau) d\tau
\end{align}

For the shifted input:
\begin{align}
y'(t) &= \int_0^t e^{-(t-\tau)} u(\tau - t_0) d\tau
\end{align}

Substitute $\sigma = \tau - t_0$:
\begin{align}
y'(t) &= \int_{-t_0}^{t-t_0} e^{-(t-\sigma-t_0)} u(\sigma) d\sigma
\end{align}

For $u(\sigma) = 0$ when $\sigma < 0$:
\begin{align}
y'(t) &= \int_0^{t-t_0} e^{-(t-t_0-\sigma)} u(\sigma) d\sigma = y(t-t_0)
\end{align}

Therefore, the system is \textbf{time-invariant} ✓

\textbf{Step 2: Convolution form}

The impulse response is $h(t) = e^{-t}u_s(t)$ where $u_s(t)$ is the unit step.

The output is:
\begin{align}
y(t) = \int_{-\infty}^{\infty} h(t-\tau)u(\tau) d\tau = \int_0^t e^{-(t-\tau)} u(\tau) d\tau = h(t) * u(t)
\end{align}

This is the convolution integral, fundamental to LTI system theory.

\textbf{Step 3: Specific example - unit step input}

For $u(t) = 1$ (unit step), $t \geq 0$:
\begin{align}
y(t) = \int_0^t e^{-(t-\tau)} \cdot 1\, d\tau = e^{-t} \int_0^t e^{\tau} d\tau = e^{-t}[e^{\tau}]_0^t = e^{-t}(e^t - 1) = 1 - e^{-t}
\end{align}

For shifted input $u(t - 1)$ (unit step at $t = 1$):
\begin{align}
y'(t) = \begin{cases}
0 & t < 1 \\
\int_1^t e^{-(t-\tau)} d\tau & t \geq 1
\end{cases}
\end{align}

For $t \geq 1$:
\begin{align}
y'(t) = e^{-t}\int_1^t e^{\tau} d\tau = e^{-t}(e^t - e) = 1 - e^{-(t-1)}
\end{align}

The shifted output is:
\begin{align}
y(t-1) = 1 - e^{-(t-1)} \quad \text{for } t \geq 1
\end{align}

Indeed, $y'(t) = y(t-1)$ ✓ Time invariance verified!

\textbf{Key Insights:}
\begin{itemize}
\item Time invariance means system parameters don't change over time
\item LTI systems can be completely characterized by their impulse response $h(t)$
\item The output is the convolution of input with impulse response
\item Time invariance simplifies analysis—we can study the system at any time
\item Non-time-invariant systems (e.g., systems with time-varying gains) require more complex analysis
\end{itemize}
\end{document}
